\documentclass[11pt,a4paper]{article}
\usepackage{../pelm}
\begin{document}
\paper{Midterm Examination}{Weeks 1--7 \quad$\cdot$\quad 30 questions $\times$ 2 points = 60 points, scaled to 40\% of the course grade}{30}{60 minutes}

\begin{questions}

% ---------------------------------------------------------------- Week 1
\q{Which of the following tasks is \textbf{generative}?}
  {Ranking job applications by predicted suitability}
  {Flagging a transaction as possibly fraudulent}
  {Drafting a discharge letter for a patient's family}
  {Detecting whether an image contains a tumour}

\q{Which statement about the term \textbf{prompt} is correct?}
  {It refers only to the text you type}
  {It includes attached documents and images}
  {It excludes any instruction added by the product}
  {It refers to the model's response}

\q{A model gives a confident, well-structured answer about a topic on which it has very little
   training material. What happens to its expressed confidence?}
  {It falls in proportion to the thinness of the coverage}
  {It stays essentially the same}
  {It rises, because uncertainty makes output longer}
  {It is reported as a percentage}

\q{An AI summary of a report is entirely accurate but omits the single most significant finding.
   Which failure signature is this?}
  {Fabrication}
  {Displacement}
  {Omission}
  {Silent failure}

\q{Why is the dangerous AI failure described in this course as ``the fluent one''?}
  {Fluent answers take longer to produce}
  {Well-formatted, confident output resembles competent professional work and is trusted accordingly}
  {Fluency indicates the model is guessing}
  {Fluent text uses more tokens}

% ---------------------------------------------------------------- Week 2
\q{ELIZA (1966) contained:}
  {A small knowledge base of medical facts}
  {A simple neural network}
  {Keyword matching and sentence reflection, with no knowledge or memory}
  {A statistical model of English grammar}

\q{What is the ``ELIZA effect''?}
  {The tendency of early chatbots to crash under load}
  {The attribution of understanding to a system that produces plausible language}
  {The improvement of output when a persona is specified}
  {The degradation of models trained on their own output}

\q{The AI winters were primarily caused by:}
  {Technical failure of the underlying methods}
  {The gap between what was promised and what was delivered}
  {Shortages of computing hardware}
  {Legal restrictions on research}

\q{Which best describes the constant direction of change across the four eras of the field?}
  {From general systems towards specialised ones}
  {From open research towards commercial secrecy}
  {From human-written rules towards learning from data}
  {From text towards images}

% ---------------------------------------------------------------- Week 3
\q{Which is \textbf{not} typically a single token?}
  {the}
  {hospital}
  {a rarely-used technical term}
  {and}

\q{A long agglutinated Turkish word is likely to be split into several tokens because:}
  {Turkish grammar is irregular}
  {The complete assembled word appeared rarely in training text, though its pieces are common}
  {Turkish uses diacritics}
  {Tokenizers cannot process non-English alphabets}

\q{Which is \textbf{not} a consequence of Turkish costing more tokens than English?}
  {The context window fills faster}
  {Metered usage costs more for the same content}
  {The model refuses Turkish input}
  {Performance is often weaker due to thinner training coverage}

\q{The model's core operation is best described as:}
  {Looking up the most relevant stored fact}
  {Producing a probability distribution over possible next tokens}
  {Applying a decision tree written by its developers}
  {Searching a database of documents}

\q{Which statement about the model's internal state is correct?}
  {It knows when it is guessing but does not report it}
  {It performs the same operation whether the pattern is strong or weak}
  {It has a separate mode for uncertain answers}
  {It calculates a confidence score for each answer}

\q{Why can these systems not simply reply ``I don't know'' whenever appropriate?}
  {The phrase is filtered out by safety systems}
  {``I don't know'' is itself a token sequence competing with all others, and it is rare after questions in training text}
  {The model is instructed never to refuse}
  {It would use too many tokens}

\q{The context window contains all of the following \textbf{except}:}
  {Your current message}
  {Documents you attached}
  {Hidden instructions added by the product}
  {Everything the model learned during training}

\q{When a document exceeds the context window, the most dangerous behaviour is:}
  {The tool refuses the input with an error}
  {Part of the material is silently dropped and a confident answer is produced anyway}
  {The response is truncated visibly}
  {The model asks you to shorten the document}

% ---------------------------------------------------------------- Week 4
\q{According to the four levers, which is \textbf{not} one of them?}
  {Context}
  {Task}
  {Format}
  {Temperature}

\q{Output is accurate and complete but reads like nothing you would ever send. Which lever should
   you pull?}
  {Context}
  {Task}
  {Format}
  {Examples}

\q{Which is the strongest reason to split ``summarise, translate, and identify risks'' into three
   separate turns?}
  {It uses fewer tokens overall}
  {Each job receives fuller attention, and you can check each result before building on it}
  {The model cannot perform three tasks}
  {Translation requires a different model}

\q{A prompt instructs: ``You are a world-class expert with thirty years of experience and three
   doctorates.'' What does this most likely achieve?}
  {A measurable improvement in factual accuracy}
  {A shift in register and vocabulary, with the superlatives adding little}
  {Access to specialist training data}
  {A reduction in hallucination}

\q{Someone adds a phrase to a prompt, the output improves, and they conclude the phrase works.
   The strongest objection is:}
  {The phrase may be too long}
  {Output varies between runs, so the improvement may be unrelated to the change}
  {The model may not have read the phrase}
  {The phrase may be in the wrong language}

\q{Which instruction converts an unverifiable claim into a checkable one?}
  {``Be accurate and thorough''}
  {``Only use reliable sources''}
  {``After each claim, quote the sentence that supports it''}
  {``Think carefully before answering''}

% ---------------------------------------------------------------- Week 5
\q{In the three dimensions of register, \textbf{field} refers to:}
  {The subject area and its vocabulary density}
  {The relationship between writer and reader}
  {The channel of communication}
  {The emotional tone}

\q{Which specification is most likely to produce reproducible style?}
  {``Make it professional and friendly''}
  {``Everyday vocabulary with no unglossed medical terms; warm but not casual; written to be read aloud''}
  {``Use an appropriate tone for the audience''}
  {``Write it well''}

\q{Why does a stated fallback such as ``write NOT STATED'' reduce invented content?}
  {It reduces the model's temperature}
  {It makes the honest answer a probable continuation rather than an unlikely one}
  {It prevents the model from accessing training data}
  {It shortens the output}

\q{Which is the most effective way to specify institutional house style?}
  {A list of adjectives describing the desired tone}
  {Two or three short samples of approved writing}
  {An instruction to write ``in our usual style''}
  {A request for formal language}

% ---------------------------------------------------------------- Weeks 6-7
\q{A generated browser tool produces a subtly wrong calculation but has a clean, working
   interface. Why is this particularly dangerous?}
  {The error is intermittent and therefore hard to reproduce}
  {The error is consistent, which makes the tool appear reliable}
  {The tool will eventually stop working}
  {The interface will change without warning}

\q{Grounding a system in your own documents changes:}
  {What the model knows}
  {Where the answer comes from}
  {Whether the model checks the truth of its output}
  {The model's training data}

\q{An assistant grounded in your documents produces an answer with no supporting quotation. The
   most likely explanation is:}
  {It over-read the passage}
  {It retrieved the wrong passage}
  {It answered from memory}
  {The context window overflowed}

\end{questions}
\end{document}
